\newpage 
\subsection{Asynchronous-update experiment}\label{apdx:async_update}

In this appendix, we report the results of our experiments where  individual agents update with independent update rates asynchronously. 

\textit{Generating the update schedule.} We use the same questions, the same 32 agents, and the same random interaction graphs as the main experiment. For each question–$J$ pair, we generate one trajectory and run it for $T=8$ timesteps. Here, we divide each step into $1000$ equal cells, so one cell is $1/1000 $ of a step. In each cell, each agent updates with a probability of 0.5/1000, independently of all other agents and cells. This means each agent updates 0.5 times per step on average (4 times over the whole run), and the times between its updates are random. Reading the sampled cells in order gives, for every step, the list of who updates and in what order. In the rare case where two agents land in the same time cell, we put them in a random order, so no two agents ever update at the same moment. 

\textit{Running the dynamics.} At time $0$, every agent answers the question with an empty inbox, exactly as in the synchronous experiment. After that, the run follows the saved schedule. When an agent's turn comes, it does what an agent does in one synchronous round: it reads its current inbox (the most recent message from each neighbor), writes a message stating its current answer, sends it to its neighbors, and then samples a new answer from the model (majority vote over $K = 5$ samples, as before). Because updates happen one at a time, an agent that updates late in a step sees the new messages of agents that updated earlier in the same step. This is the key difference from the synchronous setting, where all agents read the previous round's messages.

\begin{table}[h]
\centering
\small
\begin{tabular}{l cccc cc cc}
\toprule
& \multicolumn{4}{c}{Baselines}
& \multicolumn{2}{c}{Discrete Update}
& \multicolumn{2}{c}{Continuous Update} \\
\cmidrule(lr){2-5} \cmidrule(lr){6-7} \cmidrule(lr){8-9}
& M & P & IF & MF
& Base & + Three Couplings
& Base & + Three Couplings \\
\midrule
Subjective In-d.  & 61.1 & 78.3 & 56.3 & 63.4 & 60.6 & 66.0 & 79.7 & \textbf{80.4} \\
Subjective OOD    & 58.9 & 76.9 & 52.2 & 59.8 & 57.7 & 62.8 & 77.3 & \textbf{78.5} \\
Objective In-d.   & 46.7 & 64.3 & 48.9 & 52.8 & 51.4 & 54.9 & 63.8 & \textbf{65.3} \\
Objective OOD     & 48.2 & 65.0 & 48.0 & 55.1 & 53.0 & 54.7 & 64.7 & \textbf{65.9} \\
\bottomrule
\end{tabular}
\caption{\textit{Prediction under asynchronous dynamics.} The baselines are Majority Class (M), Persistence (P), Interaction-Free (IF), and Mean-Field (MF). The table reports the rollout accuracies. The model is fitted and evaluated on the asynchronous setup described above. GPT-4o-mini is used as the language model. \label{tab:async}}
\end{table}



We observe that the continuous model outperforms the discrete model with a larger margin under this setup. 
